\section{Equation index for every seam}\label{app:seam-equations}
This appendix is generated from the same rows as \texttt{seam\_ledger.csv}. The CSV remains authoritative for taxonomy flags, source keys, resolution evidence and successor obligations.
\subsection*{S-01: Quantum gravity / ultraviolet completion}
\begingroup\small
\begin{equation*}
\Delta\mathcal L_{2\ell}\propto \epsilon^{-1}\sqrt{-g}\,R_{\mu\nu}{}^{\rho\sigma}R_{\rho\sigma}{}^{\alpha\beta}R_{\alpha\beta}{}^{\mu\nu}\neq0;\allowbreak\quad E\ll\bar M_P\text{ EFT expansion}
\end{equation*}
\endgroup
Pure Einstein gravity has a nonzero two-loop divergence as a finite-parameter perturbative QFT, while quantum GR is predictive as a Wilsonian EFT at low energy.
\subsection*{S-02: Vacuum energy / cosmological constant}
\begingroup\small
\begin{equation*}
\Lambda_{\rm eff}=\Lambda_{\rm bare}+\bar M_P^{-2}\rho_{\rm vac}^{\rm ren};\allowbreak\quad G_{\mu\nu}+\Lambda_{\rm eff}g_{\mu\nu}=\bar M_P^{-2}T_{\mu\nu}
\end{equation*}
\endgroup
A small positive effective cosmological constant fits the background expansion. Its radiative stability and relation to quantum vacuum contributions are unexplained.
\subsection*{S-03: Classical singularities and geodesic incompleteness}
\begingroup\small
\begin{equation*}
\exists\ \gamma:[0,\lambda_*)\to\mathcal X\ \text{causal geodesic with finite }\lambda_*\text{ and no extension under theorem assumptions}
\end{equation*}
\endgroup
Singularity theorems establish causal geodesic incompleteness under specified focusing, causal and global assumptions. They do not identify a physical point of infinite density or a unique completion.
\subsection*{S-04: Dark matter identity}
\begingroup\small
\begin{equation*}
H^2=\frac{8\pi G}{3}(\rho_b+\rho_r+\rho_{\rm DM}+\rho_\Lambda);\allowbreak\quad \Omega_c h^2\simeq0.120\ \text{in base }\Lambda\mathrm{CDM}
\end{equation*}
\endgroup
Multiple cosmological and astrophysical datasets require additional gravitating structure in the standard framework, but no microscopic dark-matter particle has been confirmed.
\subsection*{S-05: Dark energy microphysics and possible evolution}
\begingroup\small
\begin{equation*}
\frac{\ddot a}{a}=-\frac{4\pi G}{3}\sum_i(\rho_i+3p_i);\allowbreak\quad w_\Lambda=-1
\end{equation*}
\endgroup
A cosmological constant gives a successful minimal fit. DESI DR2 combinations show a dataset- and parametrization-dependent preference for evolving w(a), not a settled detection.
\subsection*{S-06: Neutrino masses and lepton-number structure}
\begingroup\small
\begin{equation*}
\mathcal L_5=\frac{C_5^{ij}}{2\Lambda_L}(L_i^T C\epsilon H)(H^T\epsilon L_j)+\mathrm{h.c.},\quad m_\nu=-\frac{v^2}{2\Lambda_L}C_5
\end{equation*}
\endgroup
Oscillations require nonzero mass splittings. The renormalizable SM field content gives massless neutrinos; the dimension-five operator is the minimal data-required EFT extension.
\subsection*{S-07: Baryogenesis / matter-antimatter asymmetry}
\begingroup\small
\begin{equation*}
\eta_B\equiv n_B/n_\gamma\sim6\times10^{-10};\allowbreak\quad \dot n_B+3Hn_B=\mathcal C_{B\!\!\!/}[f_i]
\end{equation*}
\endgroup
The observed baryon asymmetry is not generated in an accepted calculation by the minimal SM with its measured Higgs and standard thermal history. Sphalerons provide B+L violation but the electroweak transition is a crossover and SM CP violation is insufficient in standard scenarios.
\subsection*{S-08: Inflation or an alternative origin of primordial perturbations}
\begingroup\small
\begin{equation*}
\mathcal P_\zeta(k)=A_s(k/k_*)^{n_s-1+\cdots};\allowbreak\quad r\equiv\mathcal P_T/\mathcal P_\zeta<0.036\ (95\%\ \mathrm{CL})
\end{equation*}
\endgroup
A nearly adiabatic, nearly scale-invariant scalar spectrum is well measured. The generating field content, initial quantum state and reheating history are not uniquely selected.
\subsection*{S-09: Electroweak hierarchy / Higgs sensitivity}
\begingroup\small
\begin{equation*}
\delta m_H^2\sim\frac{\Lambda^2}{16\pi^2}\left(6\lambda+\frac94g^2+\frac34g^{\prime2}-6y_t^2\right)\quad\text{in a cutoff parametrization}
\end{equation*}
\endgroup
The measured electroweak scale is radiatively sensitive to heavy thresholds in many UV completions. The SM as a renormalized QFT is internally consistent without a naturalness postulate.
\subsection*{S-10: Strong CP}
\begingroup\small
\begin{equation*}
\bar\theta=\theta_{\rm QCD}+\arg\det(M_uM_d);\allowbreak\quad \mathcal L_\theta=\frac{g_s^2\bar\theta}{32\pi^2}G\widetilde G;\allowbreak\quad |\bar\theta|\lesssim10^{-10}
\end{equation*}
\endgroup
QCD permits CP violation, yet neutron-EDM limits require an extremely small effective angle. The smallness is unexplained in the SM.
\subsection*{S-11: Flavour, Yukawa and mixing structure}
\begingroup\small
\begin{equation*}
\mathcal L_Y=-\bar QY_dHd_R-\bar QY_u\widetilde H u_R-\bar LY_eHe_R+\mathrm{h.c.};\allowbreak\quad Y_f\ \text{are general complex matrices}
\end{equation*}
\endgroup
Masses, mixing angles and CP phases are measured inputs. Their hierarchies and textures are not predicted by the SM gauge principle.
\subsection*{S-12: Why three generations?}
\begingroup\small
\begin{equation*}
N_g=3;\allowbreak\quad \mathcal A_{\rm local}^{(1\ generation)}=0\ \text{and}\ N_{\rm SU(2)\ doublets}^{(1\ generation)}=4\equiv0\pmod2
\end{equation*}
\endgroup
Anomaly cancellation occurs generation by generation and therefore does not fix N\_g=3 in the minimal four-dimensional SM.
\subsection*{S-13: Cosmological initial state and thermodynamic arrow}
\begingroup\small
\begin{equation*}
\rho(t_0),\ \gamma_{ij}(t_0),\ K_{ij}(t_0)\ \text{satisfy constraints};\allowbreak\quad S_{\rm coarse}(t_{\rm early})\ll S_{\rm max}
\end{equation*}
\endgroup
Field equations require a state and constrained initial data. The observed low-entropy past and perturbation state are not fixed by the GR+SM action.
\subsection*{S-14: Black-hole information and semiclassical unitarity}
\begingroup\small
\begin{equation*}
S_{\rm BH}=A/(4G\hbar);\allowbreak\quad \rho_{\rm pure}\xrightarrow{\rm semiclassical\ evaporation}\rho_{\rm thermal}\ \text{if no correlations/remnant/interior channel}
\end{equation*}
\endgroup
Semiclassical reasoning creates a tension between horizon thermality, complete evaporation and unitary quantum evolution. Observations probe exteriors, not the full quantum information map.
\subsection*{S-15: Global form of the Standard Model gauge group}
\begingroup\small
\begin{equation*}
G_{\rm SM}=\big[SU(3)_C\times SU(2)_L\times U(1)_Y\big]/\Gamma,\quad \Gamma\in\{1,\mathbb Z_2,\mathbb Z_3,\mathbb Z_6\}
\end{equation*}
\endgroup
The known local spectrum fixes the Lie algebra but not uniquely the faithful global gauge group; differences occur in line operators, monopoles and allowed heavy charges.
\subsection*{S-16: Electroweak-vacuum longevity}
\begingroup\small
\begin{equation*}
\lambda_{\rm eff}(\mu)\ \text{may cross zero};\allowbreak\quad \Gamma/V\sim \mu_B^4e^{-S_B[\lambda_{\rm eff},\ldots]}
\end{equation*}
\endgroup
High-scale extrapolation often places the SM near a metastability boundary; the conclusion depends on top-mass interpretation, alpha\_s, matching and higher-dimensional operators.
\subsection*{S-17: Accidental baryon/lepton symmetries and proton stability}
\begingroup\small
\begin{equation*}
\mathcal L_{\rm SMEFT}\supset C_{qqql}\,qqql/\Lambda^2+\cdots;\allowbreak\quad B,L\ \text{are accidental at dimension }4
\end{equation*}
\endgroup
Renormalizable SM interactions conserve baryon and lepton number perturbatively (with electroweak anomaly qualifications), while generic higher-dimensional operators need suppression.
\subsection*{S-18: Active precision-anomaly watchlist}
\begingroup\small
\begin{equation*}
\Delta_O=(O_{\rm data}-O_{\rm baseline})/\sqrt{\sigma_{\rm data}^2+\sigma_{\rm theory}^2};\allowbreak\quad \Delta_O\ \text{is dataset/theory-version dependent}
\end{equation*}
\endgroup
H0, muon g-2, the CDF II W mass, selected CKM/flavour quantities and DESI dark-energy fits deserve continued scrutiny. Their status is not uniform and changes with data/theory updates.
